%%
%% CS 360 Final Project Report -- Phase IV
%%
\documentclass[sigconf,nonacm]{acmart}

\AtBeginDocument{%
  \providecommand\BibTeX{{Bib\TeX}}}

\begin{document}

\title{Bull Market: A Web-Based Sports Prediction Market
       with Insider Trader Detection}

\author{Andrew Araujo}
\email{arau2523@vandals.uidaho.edu}
\affiliation{%
  \institution{University of Idaho}
  \city{Moscow}
  \state{Idaho}
  \country{USA}
}

\author{Jason Culbertson}
\email{culb1469@vandals.uidaho.edu}
\affiliation{%
  \institution{University of Idaho}
  \city{Moscow}
  \state{Idaho}
  \country{USA}
}

\author{Xavier Dahlstet}
\email{dahl6044@vandals.uidaho.edu}
\affiliation{%
  \institution{University of Idaho}
  \city{Moscow}
  \state{Idaho}
  \country{USA}
}

\renewcommand{\shortauthors}{Araujo, Culbertson, Dahlstet.}

\begin{abstract}
We built Bull Market, a web-based sports prediction market where users
register, deposit funds, and place bets on live sports events. The
platform is built with Node.js, Express, and TypeScript on the backend,
connected to a MySQL~8.0 database with ten normalized tables. The system
integrates live market data from the Polymarket API and supports both
straight bets and parlay bets. Motivated by recent research on insider
trading in real-world prediction markets, we also developed an insider
trader detection module that applies Z-Score analysis and K-Means
clustering to identify suspicious betting patterns in a simulated
dataset of 120 users and 80 events. The completed system includes nine
web interfaces, a RESTful API backend, full financial transaction
management, data visualization with Chart.js, and a research component
with precision, recall, and F1 evaluation of both detection methods.
\end{abstract}

\keywords{Prediction Markets, Database Systems, Web Development,
Insider Trading Detection, Unsupervised Machine Learning}

\begin{teaserfigure}
  \centering
  \includegraphics[width=\textwidth, height=0.75\textwidth]{erdiagram.png}
  \caption{Entity-Relationship diagram for the Bull Market database
  schema, showing the ten tables and their relationships across users,
  accounts, events, odds, bets, and research data.}
  \label{fig:erdiagram}
\end{teaserfigure}

\maketitle

%%-----------------------------------------------------------
\section{Introduction}

Bull Market is a sports prediction market platform built as our
CS~360 database systems project. A prediction market lets users bet
money on the outcomes of future events, with odds that reflect the
collective beliefs of all participants. Our platform lets registered
users browse upcoming sports events, place bets based on current odds,
and track their winnings and losses over time. An admin panel controls
event creation, odds management, and resolution, automatically paying
out winning bets when an event is settled.

The project was also shaped by two recent papers on real-world
prediction markets (discussed in Section~\ref{sec:relatedwork}). Those
papers documented serious insider trading on Polymarket, the same
platform whose API we use on our homepage, and raised concerns about
how prediction markets can hide important information from users. These
findings motivated us to add a research module to Bull Market: an
insider trader detection system that uses two unsupervised machine
learning methods to flag users with suspiciously strong performance.

This report covers the full Phase~IV deliverables: our finalized
architecture, database design, the three core database queries,
implementation details, the complete working prototype, the insider
detection system, and an evaluation of what we built and what we would
improve next.

%%-----------------------------------------------------------
\section{Related Work}
\label{sec:relatedwork}

Two recent papers on Polymarket directly shaped our project design.

\subsection{Informed Trading in Prediction Markets}

Mitts and Ofir~\cite{mitts2025iran} studied insider trading on
Polymarket by analyzing 210,718 suspicious wallet-market pairs between
February 2024 and February 2026. The traders they flagged achieved a
69.9\% win rate, more than 60 standard deviations above random
performance, and collectively earned approximately \$143 million in
anomalous profits.

The authors documented specific incidents where large positions were
opened just before public announcements of geopolitical events,
including Iranian military operations and the Venezuelan election
result. These positions were almost certainly placed by traders with
advance knowledge of the outcome. Because prediction markets settle
on real-world facts, inside information carries a nearly certain
payoff. The paper also identifies patterns useful for detection: the
flagged traders placed an unusually high fraction of their bets very
close to the event resolution time, had highly consistent stake sizes,
and achieved win rates far above what skill alone could explain.

The paper concludes that existing U.S.\ securities law does not cleanly
cover insider trading in prediction markets, since these markets
technically deal in event contracts rather than securities. The authors
argue that the CFTC could extend oversight to cover them.

This paper directly motivated our insider detection module. The same
behavioral signals identified in that paper (high win rates, high ROI,
and bets placed close to event start times) are exactly the four
features we extract and analyze with our Z-Score and K-Means methods
(see Section~\ref{sec:research}).

\subsection{Prediction Laundering}

Rohanifar, Ahmed, and Sultana~\cite{rohanifar2025laundering} conducted
a critical sociotechnical audit of how Polymarket curates markets and
presents information to users. They argue the platform creates a false
sense of neutrality through four stages they call ``prediction
laundering'': structural sanitization (selecting which future events
can be bet on), probabilistic flattening (reducing complex, contested
events to a single percentage), architectural masking (concealing the
influence of large whale traders on market prices), and epistemic
hardening (erasing governance disputes so the final outcome appears
clean and factual).

A key finding is that market prices on Polymarket are heavily shaped
by a small number of high-capital traders, but the interface presents
probabilities as though they represent the aggregate wisdom of many
independent participants. Retail users have no way to see how much
of the price is driven by a single actor. The authors call for
``friction-positive'' interface design, meaning systems that expose
the social and financial dynamics driving a market rather than smoothing
them over.

This paper influenced our design philosophy. Rather than hiding how
our platform works, we built the admin and research pages to be
deliberately transparent. The admin page shows exactly how events are
created, how odds are set, and how events are resolved. The research
page exposes every feature used in detection and lets users compare
both methods side by side before and after revealing the ground truth.
We also display raw decimal odds on the Event page rather than
converting them to a single probability, giving users more information
about how the odds are structured.

%%-----------------------------------------------------------
\section{Technology Stack}

\subsection{Frontend}
All nine pages are standalone HTML files using HTML5, CSS3, and vanilla
JavaScript. We used Bootstrap~5 for responsive layout and styling, and
Chart.js for data visualization. Chart.js drives the seven-day odds
history line chart on the Event page and the Win Rate vs.\ ROI scatter
plot on the Research page.

\subsection{Backend}
The backend is a Node.js server running Express~5, written in
TypeScript and executed with \texttt{tsx}. All business logic and SQL
queries live in \texttt{server.ts}, which exposes roughly 40 REST
endpoints. A separate \texttt{db.ts} file manages the MySQL connection
pool using the \texttt{mysql2} library, reading credentials from a
\texttt{.env} file that is excluded from version control.

\subsection{Database}
We use MySQL~8.0 managed locally through XAMPP. The schema has ten
tables: eight core business tables plus two supporting tables for odds
history and research ground truth. All ten tables are created
automatically at server startup using \texttt{CREATE TABLE IF NOT
EXISTS}, so the application self-initializes on first run.

\subsection{Other Tools}
Password hashing uses the \texttt{bcrypt} library with a salt round
factor of 10. Client-side sessions are stored in \texttt{localStorage}.
The Polymarket Gamma API supplies live market data for the homepage.
The \texttt{@faker-js/faker} library generates synthetic users, events,
and betting records for the research module.

Our original plan called for React, Prisma, Socket.io, and AWS hosting.
We pivoted to vanilla HTML/JS because it was faster to iterate, and
XAMPP simplified local database management. React's build toolchain
added overhead that slowed development, and Prisma's abstraction made
it harder to write and debug the custom SQL needed for the betting
transaction and research queries. TypeScript was kept throughout the
backend for type safety.

%%-----------------------------------------------------------
\section{Project Plan}

We followed a seven-week sprint methodology, treating each week as an
Agile sprint with assigned tasks and a weekly goal. At the start of
each sprint we identified the highest-priority features, divided
responsibilities across the three team members, and reviewed progress
from the previous week. The team used a GitHub project board to track
open tickets and pull requests throughout development.

The most significant unplanned change was the addition of the Research
module. After reviewing the two prediction market papers midway through
development, we decided the insider detection component would be a
strong complement to the core betting platform. This required a new
database table (\texttt{research\_seed\_truth}), new API routes, a
complete seeding system, and the Research page UI. Most of this work
landed in Weeks~5 and~6 rather than following the original plan.

We also pivoted away from React and Prisma earlier than expected.
By the end of Week~1 it was clear that the React build system was
slowing us down, and vanilla HTML/JS let us move faster while still
meeting all the database-focused requirements of the course.

%%-----------------------------------------------------------
\section{Timeline of Deliverables}

\subsection{Week 1: Foundation and Authentication}
\textit{Goal: Register and login working.} We built the Express
server structure, defined the initial MySQL schema, and implemented
user registration and sign-in with bcrypt password hashing. The
database connection pool in \texttt{db.ts} was established and tested.

\subsection{Week 2: Database and Schema}
\textit{Goal: All core tables live, deposit working.} We finalized
all eight core tables with foreign key constraints and implemented the
account deposit and withdrawal endpoints. The Accounts and Transactions
tables were connected end-to-end through the Account page.

\subsection{Week 3: Events and Odds}
\textit{Goal: Markets browsable, live Polymarket data displaying.} We
built the Events and Odds tables, implemented the markets browsing
page with sport and status filters, and integrated the Polymarket Gamma
API to display live markets on the homepage.

\subsection{Week 4: Bet Placement}
\textit{Goal: Atomic bets, balance correctly deducted.} We implemented
the bet placement transaction with the \texttt{FOR UPDATE} row lock,
added parlay support via the BetSelections table, and built the Event
detail page with the bet slip and real-time payout calculation.

\subsection{Week 5: Payouts and Research Start}
\textit{Goal: Full money loop complete, research module started.} We
implemented event resolution and automatic bet settlement for all
associated bets, completed the Dashboard and Account pages, and began
building the insider detection module after reviewing the two papers.

\subsection{Week 6: Frontend and Research Completion}
\textit{Goal: All interfaces complete, research module working.} We
completed the Research page with both detection methods, the scatter
plot visualization, the ground truth reveal, and precision, recall, and
F1 metrics. The \texttt{odds\_history} table was added to support the
Chart.js line chart on the Event page.

\subsection{Week 7: Testing and Polish}
\textit{Goal: Demo-ready across all nine pages.} We seeded complete
test datasets, fixed edge cases in bet settlement, finalized the admin
controls, and cleaned up the UI. Final end-to-end testing confirmed
the full money loop works correctly.

%%-----------------------------------------------------------
\section{Finalized System Architecture}

The system has three layers that communicate over HTTP.

\textbf{Frontend (nine HTML pages):} Each page is self-contained and
communicates with the backend through JavaScript \texttt{fetch()} calls.
Pages receive JSON responses and render data dynamically without a
full-page reload.

\textbf{Backend (Express API):} \texttt{server.ts} implements all REST
endpoints, organized into six logical groups: authentication
(\texttt{/register}, \texttt{/login}), accounts
(\texttt{/api/account}), events and odds (\texttt{/api/events}), bets
(\texttt{/api/bets}), admin (\texttt{/api/admin}), and research
(\texttt{/api/research}).

\textbf{Database (MySQL~8.0):} Ten tables with foreign key constraints
enforce data integrity. The typical data flow is: a page sends a
\texttt{fetch()} request to an Express route, the route executes a SQL
query through the connection pool, and the result is returned as JSON
for the page to render.

\textbf{Visualization:} Two data visualizations are built with
Chart.js. The Event page renders a multi-line chart of decimal odds
over the past seven days, with one line per selection (e.g., Yes and
No). The Research page renders a scatter plot of Win Rate vs.\ Average
ROI, colored by the active detection method's suspect labels or cluster
assignments. Both charts update dynamically when new data arrives from
the API.

%%-----------------------------------------------------------
\section{Database Design}

\subsection{Schema}

\noindent
Sports(\underline{sport\_id}, name)\\
Users(\underline{user\_id}, username, email, password\_hash,
  date\_of\_birth, created\_at)\\
Accounts(\underline{account\_id}, user\_id, balance, currency)\\
Transactions(\underline{transaction\_id}, account\_id, user\_id, type,
  amount, status, created\_at)\\
Events(\underline{event\_id}, sport\_id, name, start\_time, status,
  result, description)\\
Odds(\underline{odds\_id}, event\_id, market\_type, selection\_name,
  decimal\_odds, updated\_at)\\
Bets(\underline{bet\_id}, user\_id, stake\_amount, total\_odds,
  potential\_payout, status, placed\_at)\\
BetSelections(\underline{selection\_id}, bet\_id, odds\_id, event\_id,
  odds\_at\_placement, result, selection\_type)\\
OddsHistory(\underline{history\_id}, odds\_id, event\_id, decimal\_odds,
  recorded\_at)\\
ResearchSeedTruth(\underline{id}, user\_id, is\_insider, archetype)\\

\begin{figure}[h!]
  \centering
  \includegraphics[width=\linewidth]{databasestate.png}
  \caption{The MySQL users table viewed in phpMyAdmin, showing
  registered users with auto-generated usernames and bcrypt-hashed
  passwords.}
  \label{fig:dbstate}
\end{figure}

\subsection{Normalization}

\subsubsection{First Normal Form}
First normal form requires a primary key, no arrays stored in a single
cell, and no repeating groups of columns~\cite{gomstyn2024dbnorm}.
Every table has a single-column surrogate primary key. Parlay bets
store each individual pick as a separate row in BetSelections rather
than as a list inside the Bets table. OddsHistory stores each recorded
price as its own row rather than as an array. All columns hold a single
atomic value.

\subsubsection{Second Normal Form}
Second normal form requires that every non-key attribute depends on the
whole primary key, not just part of it, which is only relevant for tables with
composite keys~\cite{gomstyn2024dbnorm}. All of our tables use single
surrogate keys (bet\_id, selection\_id, etc.), so second normal form is
automatically satisfied with no changes required.

\subsubsection{Third Normal Form}
Third normal form forbids transitive dependencies, where a non-key
column depends on another non-key column instead of the primary
key~\cite{gomstyn2024dbnorm}. One deliberate exception is
\texttt{potential\_payout} in Bets, which is derivable from
\texttt{stake\_amount}~$\times$~\texttt{total\_odds}. We store it as
a snapshot because odds change in real time, and once a bet is placed
the agreed payout must be frozen permanently. The same reasoning
applies to \texttt{odds\_at\_placement} in BetSelections, which
intentionally diverges from the live \texttt{decimal\_odds} in the
Odds table after it is recorded.

\subsubsection{BCNF}
Boyce-Codd Normal Form requires that the left-hand side of every
functional dependency is a superkey~\cite{gomstyn2024dbnorm}. All
attributes in our schema depend only on their own table's primary key,
with no hidden candidate keys that would cause a violation.

%%-----------------------------------------------------------
\section{Required Queries}

Three SQL queries form the core of the platform's functionality.

\subsection{Query 1: Atomic Bet Placement}
Placing a bet requires multiple writes that must all succeed or all
fail together. We use a MySQL transaction with a
\texttt{SELECT~...~FOR~UPDATE} row lock to prevent race conditions
where two concurrent requests could both read the same balance and
overdraft the account.

\begin{verbatim}
START TRANSACTION;

SELECT account_id, balance
FROM accounts
WHERE user_id = ? FOR UPDATE;

UPDATE accounts
  SET balance = balance - ?
  WHERE user_id = ?;

INSERT INTO bets
  (user_id, stake_amount,
   total_odds, potential_payout)
  VALUES (?, ?, ?, ?);

INSERT INTO bet_selections
  (bet_id, odds_id, event_id,
   odds_at_placement, selection_type)
  VALUES (?, ?, ?, ?, ?);

INSERT INTO transactions
  (account_id, user_id, type,
   amount, status)
  VALUES (?, ?, 'bet', ?, 'completed');

COMMIT;
\end{verbatim}

\subsection{Query 2: Odds History for Visualization}
The Event page shows a Chart.js line chart of how odds moved over
the past seven days. This query pulls the full recorded history for
all selections in a given event, ordered so Chart.js can group each
line by selection name.

\begin{verbatim}
SELECT oh.recorded_at,
       oh.decimal_odds,
       o.selection_name,
       o.market_type,
       o.odds_id
FROM odds_history oh
JOIN odds o
  ON oh.odds_id = o.odds_id
WHERE oh.event_id = ?
ORDER BY o.selection_name,
         oh.recorded_at ASC;
\end{verbatim}

\subsection{Query 3: Insider Detection Feature Extraction}
The research analysis begins by pulling all settled bet data in a
single query joining four tables. The four detection features are then
computed in application code from this result set. The
\texttt{TIMESTAMPDIFF} call gives the number of minutes between bet
placement and event start, which is used to flag late bets.

\begin{verbatim}
SELECT u.user_id, u.username,
       b.stake_amount,
       b.potential_payout,
       b.status AS bet_status,
       TIMESTAMPDIFF(MINUTE,
         b.placed_at,
         e.start_time) AS mins_before
FROM users u
JOIN bets b
  ON b.user_id = u.user_id
JOIN bet_selections bs
  ON bs.bet_id = b.bet_id
JOIN events e
  ON e.event_id = bs.event_id
WHERE b.status IN ('won', 'lost');
\end{verbatim}

%%-----------------------------------------------------------
\section{Implementation Details}

All nine pages communicate with the backend through a consistent
pattern: the page's JavaScript sends a \texttt{fetch()} request to an
Express route in \texttt{server.ts}, the route executes a SQL query
through the connection pool in \texttt{db.ts}, and the result is
returned as JSON. Credentials are stored in a \texttt{.env} file and
excluded from version control so the database password is never
committed to the repository.

\subsection{User Registration}
The registration endpoint hashes the submitted password with bcrypt
before writing anything to the database, so plaintext passwords are
never stored. It also auto-generates a username from the user's first
and last name and creates a linked account row with a zero balance in a
separate \texttt{INSERT}.

\begin{verbatim}
app.post('/register',
  async (req, res) => {
  const { firstname, lastname,
          email, password } = req.body;

  const hash =
    await bcrypt.hash(password, 10);
  const username =
    `${firstname}${lastname}`
      .toLowerCase();

  const [result] = await db.query(
    `INSERT INTO users
      (username, email,
       password_hash, date_of_birth)
     VALUES (?, ?, ?, ?)`,
    [username, email, hash,
     '2000-01-01']
  );
  const { insertId } = result;

  await db.query(
    `INSERT INTO accounts
      (user_id) VALUES (?)`,
    [insertId]
  );
  res.status(201).json(
    { success: true }
  );
});
\end{verbatim}

\subsection{User Authentication}
The sign-in endpoint accepts either an email address or a username
by checking for the presence of an \texttt{@} symbol in the submitted
value. Bcrypt's \texttt{compare} function checks the submitted password
against the stored hash without ever decrypting it.

\begin{verbatim}
app.post('/login',
  async (req, res) => {
  const { email, password } = req.body;
  const isEmail = email.includes('@');

  const [rows] = await db.query(
    isEmail
      ? `SELECT * FROM users
           WHERE email = ?`
      : `SELECT * FROM users
           WHERE username = ?`,
    [email]
  );
  const user = rows[0];
  if (!user) {
    res.status(401).json(
      { error: 'Invalid credentials' }
    );
    return;
  }
  const ok = await bcrypt.compare(
    password, user.password_hash
  );
  if (!ok) {
    res.status(401).json(
      { error: 'Invalid credentials' }
    );
    return;
  }
  res.json({
    success: true,
    user: {
      id: user.user_id,
      username: user.username,
      email: user.email
    }
  });
});
\end{verbatim}

%%-----------------------------------------------------------
\section{Complete Working Prototype}

The system includes nine fully functional interfaces.

\subsection{Home}
The landing page fetches live prediction markets from the Polymarket
Gamma API and displays them alongside our own internal events. Each
market shows the current probability as a percentage derived from live
trading activity on Polymarket.

\subsection{Markets}
Users can browse all internal events, filter by sport category, and
filter by status (upcoming, live, or finished). Clicking an event
navigates to its detail page.

\subsection{Event and Visualization}
The event detail page shows current decimal odds for each selection and
a seven-day odds history chart rendered with Chart.js. As the admin
updates odds over time, the history table records each change, and the
chart displays all selections as separate colored lines. Users can enter
a stake amount and the page calculates the potential payout in real
time. Submitting places the bet if the user's balance is sufficient.

\subsection{Sign In and Register}
The Sign In page accepts either an email address or a username and
authenticates using bcrypt comparison. The Register page collects a
first name, last name, email, and password, auto-generates a username,
and creates the account with a zero balance.

\subsection{Dashboard}
Shows the user's recent bets with stake amount, potential payout, and
current status. Resolved bets are color-coded: green for wins and red
for losses, giving users a quick visual summary of their P\&L.

\subsection{Account}
Displays the user's current balance and allows them to deposit or
withdraw funds. The page also shows the last 50 transactions so users
can review their complete history.

\subsection{Administration}
The admin page provides tools for platform management: creating events
and odds, resolving events (which automatically settles all associated
bets and pays winners), seeding sports and research data, adjusting all
account balances, and clearing the database. It is intended for
demonstration and grading use.

\subsection{Research}
The Research page runs insider trader detection on the seeded dataset
(see Section~\ref{sec:research}). Users click ``Run Both Methods'' to
execute the analysis, then switch between tabs to compare Z-Score and
K-Means results. A scatter plot updates to reflect the active tab.
Clicking ``Reveal Ground Truth'' fetches labels from the
\texttt{research\_seed\_truth} table and displays precision, recall,
and F1 scores for each method.

\begin{figure}[h!]
  \centering
  \includegraphics[width=\linewidth]{home.png}
  \caption{The Bull Market homepage displaying live Polymarket markets
  and internal events with current outcome probabilities.}
  \label{fig:home}
\end{figure}

\begin{figure}[h!]
  \centering
  \includegraphics[width=\linewidth]{signin.png}
  \caption{The Sign In page, which accepts either an email address or
  a username and authenticates with bcrypt.}
  \label{fig:signin}
\end{figure}

\begin{figure}[h!]
  \centering
  \includegraphics[width=\linewidth]{register.png}
  \caption{The Register page for creating a new account. A username is
  auto-generated from the user's first and last name.}
  \label{fig:register}
\end{figure}

%%-----------------------------------------------------------
\section{Insider Trader Detection}
\label{sec:research}

Motivated by the findings of Mitts and
Ofir~\cite{mitts2025iran}, we built a research module that applies
two unsupervised machine learning methods to detect insider traders
in a controlled, simulated dataset. Neither method is given any labels
during training; they discover suspicious users entirely from
behavioral patterns in the betting data.

\subsection{Simulated Dataset}

The dataset is generated using \texttt{@faker-js/faker} and seeded
from the Admin page. It consists of 120 users across five behavioral
archetypes, 80 events, and hundreds of settled bets:

\begin{itemize}
  \item \textbf{Insider (10 users):} High win rates, large stakes
  placed within 4 hours of event start. Includes deliberate losses
  to avoid a suspicious 100\% win rate and better simulate reality.
  \item \textbf{Sharp Normal (30 users):} Skilled bettors with
  above-average win rates through research and judgment, not inside
  information.
  \item \textbf{Average (30 users):} Typical recreational bettors
  with moderate and mixed results.
  \item \textbf{Recreational (30 users):} Casual bettors with small,
  inconsistent stakes and below-average win rates.
  \item \textbf{High Roller (20 users):} Large stakes but essentially
  random outcomes; they spend a lot but have no edge.
\end{itemize}

The 10 insider users are the ground truth for evaluating detection
accuracy. Their labels are stored in the
\texttt{research\_seed\_truth} table, hidden from the detection
algorithms, and only revealed when the user clicks the Reveal button.
Users must have at least ten settled bets before they are included in
the analysis to ensure statistical reliability.

\subsection{Detection Features}

Both methods use the same four features per user, grounded in the
behavioral signals identified by Mitts and Ofir~\cite{mitts2025iran}:

\begin{enumerate}
  \item \textbf{Win Rate:} Wins divided by total settled bets.
  \item \textbf{Average ROI:} Average net profit per bet relative to
  the stake. A winning bet at 2.0x decimal odds yields $+100\%$ ROI;
  a losing bet yields $-100\%$ ROI.
  \item \textbf{Late Bet Ratio:} The fraction of bets placed within
  4 hours of the event's start time. This captures the pattern of
  acting on last-minute, non-public information.
  \item \textbf{Stake Consistency:} The inverted coefficient of
  variation of a user's stake amounts ($1 - \sigma/\mu$). Insiders
  tend to bet large, consistent amounts when they have information,
  while recreational bettors have high variability in their stakes.
\end{enumerate}

\subsection{Z-Score Method}

Each feature is standardized across all users using the Z-score
formula: $z = (x - \mu) / \sigma$. The four Z-scores are summed into
a composite suspicion score. Stake consistency is inverted before
scoring (low variability should add to suspicion, not subtract from
it). Users are ranked from highest to lowest composite score, and the
top 15 are flagged as suspected insiders. This is a heuristic
method; it measures relative extremity but does not learn any
underlying structure in the data.

\subsection{K-Means Method}

K-Means clustering ($k=2$) is applied to min-max normalized feature
values using k-means++ initialization. The k-means++ strategy selects
the first centroid at random and then probabilistically selects each
subsequent centroid with likelihood proportional to its squared
distance from the nearest existing centroid, which avoids the poor
initializations that can occur with random seeding. The algorithm then
iterates assignment and update steps until centroids converge.

After clustering, the cluster with the higher mean win rate is labeled
the ``suspected insiders'' cluster. Unlike Z-Score, K-Means does not
require a fixed threshold; the cluster size is determined by the
data's natural structure.

\begin{table}[h]
  \caption{Comparison of the two detection methods.}
  \label{tab:methods}
  \renewcommand{\arraystretch}{1.2}
  \begin{tabular}{|l|l|l|}
    \hline
    \textbf{Property} & \textbf{Z-Score} & \textbf{K-Means} \\
    \hline
    Type              & Heuristic ranking & Clustering \\
    Normalization     & Z-score ($\mu$=0) & Min-max (0--1) \\
    Output            & Suspicion score   & Cluster label \\
    Threshold         & Fixed top-15      & Data-driven \\
    Initialization    & N/A               & k-means++ \\
    Interpretability  & High              & Medium \\
    \hline
  \end{tabular}
\end{table}

\subsection{Visualization and Evaluation}

After both methods run, the Research page displays a ranked table for
each method and a scatter plot of Win Rate vs.\ Average ROI, colored
by suspect status or cluster assignment. The scatter plot makes it
easy to see visually how well each method separates the insider cluster
from the general population. Clicking ``Reveal Ground Truth'' fetches
the true labels from \texttt{research\_seed\_truth} and computes
precision, recall, and F1 score for each method, making it easy to
compare their relative effectiveness.

%%-----------------------------------------------------------
\section{Evaluation and Limitations}

\subsection{What Works}
The platform functions as a complete prediction market from end to end.
Users can register, deposit, browse live events, place straight and
parlay bets, and track their P\&L on the dashboard. Winning bets are
automatically settled with the correct payout when an admin resolves
an event. The atomic transaction for bet placement correctly handles
concurrent requests without overdrafting accounts. The research module
runs both detection methods on demand and produces quantitative
accuracy metrics, making it straightforward to compare their
performance.

The integration of the Polymarket API demonstrates that the platform
could be extended to display data from external prediction markets.
The odds history visualization shows a full week of price movement and
reflects changes immediately when new odds are entered through the
admin panel.

\subsection{Limitations}
\begin{itemize}
  \item \textbf{Simulated money:} The platform is not connected to any
  real payment processor. All account balances are virtual and exist
  only in the database.
  \item \textbf{Local only:} The system requires XAMPP running locally
  and cannot be accessed by external users without additional server
  configuration.
  \item \textbf{Session security:} User sessions are stored in
  \texttt{localStorage}, which is not appropriate for a production
  environment. A real deployment would use JWTs or server-side
  session management with secure HTTP-only cookies.
  \item \textbf{Synthetic research data:} The insider detection module
  is evaluated on synthetic data only. Real-world detection performance
  would likely differ since real insiders can adapt their behavior to
  avoid detection.
  \item \textbf{No email verification:} Users can register with any
  email address without confirming ownership.
  \item \textbf{Polymarket display only:} Live Polymarket data appears
  on the homepage for context, but users can only bet on events created
  through our admin panel.
\end{itemize}

%%-----------------------------------------------------------
\section{Future Improvements}

Given more time, we would prioritize the following improvements:

\begin{itemize}
  \item \textbf{Real payments:} Integrate a payment processor such as
  Stripe for actual deposits and withdrawals, including KYC checks
  required for regulated gambling platforms.
  \item \textbf{Cloud hosting:} Deploy to AWS, Railway, or a similar
  platform so the site is publicly accessible without requiring XAMPP.
  \item \textbf{Secure sessions:} Replace \texttt{localStorage} with
  server-side session management or signed JWT tokens stored in
  HTTP-only cookies.
  \item \textbf{Real-time updates:} Add WebSocket support using
  Socket.io so odds update live across all connected clients without
  requiring a page refresh.
  \item \textbf{Live insider detection:} Run detection algorithms
  continuously on real user bets as they accumulate, rather than only
  on seeded synthetic data.
  \item \textbf{Better ML models:} Explore more sophisticated
  approaches such as isolation forests, random forests, or a simple
  neural network classifier for insider trader detection.
  \item \textbf{Mobile UI:} Improve the responsive design so all pages
  work well on smaller screens.
  \item \textbf{Bet limits and responsible gambling:} Add daily deposit
  limits, self-exclusion options, and loss-limit warnings, which would
  be legally required on a real platform.
\end{itemize}

%%-----------------------------------------------------------
\nocite{jamilCS360}
\bibliographystyle{ACM-Reference-Format}
\bibliography{sample-base}

\end{document}
\endinput
